\documentclass[11pt]{article}

% The official *ACL style (ACL, EMNLP, NAACL, EACL, ARR).
% [review] for submission, [final] for the camera-ready version.
\usepackage[review]{acl}

\usepackage{times}
\usepackage{latexsym}
\usepackage[T1]{fontenc}
\usepackage[utf8]{inputenc}
\usepackage{microtype}
\usepackage{inconsolata}
\usepackage{graphicx}
\usepackage{amsmath}

\title{{{TITLE}}}

\author{First Author \\
  University One \\
  \texttt{first@example.com} \\\And
  Second Author \\
  University Two \\
  \texttt{second@example.com} \\}

\begin{document}
\maketitle

\begin{abstract}
One paragraph: the problem, what you did, what you found, and why it matters. Keep it to about 150 words and avoid citations here.
\end{abstract}

\section{Introduction}
\label{sec:intro}

Start with the problem and why it matters. Say who is affected, what goes
wrong today, and why existing approaches fall short. Prior work such as
\cite{mcmahan2017communication} trains models without collecting data in one
place, and attention-based models~\cite{vaswani2017attention} changed how we
build them.

This paper makes three contributions:
\begin{itemize}
  \item a clear statement of the problem and the setting we study;
  \item a method that addresses it (Section~\ref{sec:method});
  \item an evaluation showing when and why it works.
\end{itemize}

\section{Method}
\label{sec:method}

Describe your approach so that a reader could reproduce it. Equations are
numbered automatically:
\begin{equation}
  \mathcal{L}(\theta) = -\sum_{i=1}^{N} \log p_\theta(y_i \mid x_i).
  \label{eq:loss}
\end{equation}
We minimise Equation~\eqref{eq:loss} with stochastic gradient descent.

\section{Evaluation}
\label{sec:eval}

Explain the setup, the baselines, and what the results show.

\section{Conclusion}

Summarise what you found and what remains open.


\section*{Limitations}

*ACL venues require a limitations section. It does not count towards the page limit.

\bibliography{references}

\end{document}
